
This work reports an empirical finding that contradicts prevailing assumptions about uncertainty quantification methods for hallucination detection: semantic entropy, which captures meaning-level uncertainty through NLI-based clustering, underperforms a simple confidence score by 0.24 AUROC points on multiple-choice hallucination detection.

Automatic detection of hallucinations---confident-sounding but factually incorrect LLM outputs---is a prerequisite for safe deployment in high-stakes domains. Uncertainty quantification offers a principled approach: outputs with high uncertainty may be unreliable and warrant verification. Recent work has produced increasingly sophisticated UQ methods, from token-level entropy and confidence measures to semantic entropy with NLI-based clustering \citep{kuhn2023semantic}, self-consistency checks \citep{manakul2023selfcheckgpt}, and P(True) calibration \citep{kadavath2022language}.

A methodological problem complicates method selection: each approach has been evaluated under favorable conditions using different benchmarks, data splits, and evaluation protocols. Kuhn et al. demonstrated semantic entropy's superiority over token entropy on custom QA splits. Manakul et al. showed SelfCheckGPT's effectiveness on WikiBio generation. Kadavath et al. validated P(True) on proprietary data. No controlled head-to-head comparison exists on identical conditions, leaving practitioners without guidance for method selection.

This work addresses this gap through a controlled comparison of UQ methods for hallucination detection on TruthfulQA's multiple-choice format with Llama-3-8B-Instruct. The central finding is that \textbf{UQ method effectiveness depends on task format}: token-level methods extract discriminative signals directly from logits without format dependency, while semantic entropy's NLI-based clustering requires substantial semantic content that single-letter multiple-choice answers cannot provide. The clustering mechanism functions correctly---verification confirms that NLI models produce clusters with non-trivial entropy variance---but the resulting entropy scores do not correlate with hallucination status on this format.

The contributions are: (1) empirical demonstration that token-level UQ methods (max probability, choice entropy) achieve AUROC 0.77--0.81 on TruthfulQA mc1 with Llama-3-8B-Instruct, exceeding a 0.55 chance threshold; (2) documentation that semantic entropy achieves only AUROC 0.5645 on the same task despite correct mechanism implementation; (3) identification of format dependency as the root cause---semantic entropy requires free-form responses with semantic content suitable for NLI comparison, while multiple-choice answers are single letters with no such content.

